% Generated draft: requires mathematical and source review.
\documentclass[11pt]{article}
\usepackage[margin=25mm]{geometry}
\usepackage{fontspec}
\setmainfont{Latin Modern Roman}
\usepackage{amsmath,amssymb,mathtools,unicode-math}
\setmathfont{Latin Modern Math}
\usepackage{enumitem,needspace,xurl,hyperref}
\setlength{\parindent}{0pt}
\setlength{\parskip}{5pt}
\setlength{\emergencystretch}{3em}
\newcommand{\N}{\mathbb{N}}
\newcommand{\Z}{\mathbb{Z}}
\newcommand{\Q}{\mathbb{Q}}
\newcommand{\R}{\mathbb{R}}
\newcommand{\C}{\mathbb{C}}
\newcommand{\F}{\mathbb{F}}
\newcommand{\A}{\mathbb{A}}
\newcommand{\PP}{\mathbb{P}}
\newcommand{\E}{\mathbb{E}}
\newcommand{\eps}{\varepsilon}
\newcommand{\dd}{\,\mathrm{d}}
\begin{document}
\section*{UnsolvedMath ID: 801059}
\textbf{Source problem code:} catalogue entry
\subsection{Definitions and mathematical statement}
The supplied record names a family of maps \(f_{c,b}\), but does not specify its formula, phase space, or admissible parameter values. It also does not define the quantities \(\lambda^-(p)\) and \(\lambda^+(p)\). These omissions remain unresolved below.

In the usual setting of a diffeomorphism, a periodic saddle is a hyperbolic periodic point with both stable and unstable directions. For such a point \(p\), let \(\mathcal O(p)\) denote its periodic orbit, and let \(W^s(\mathcal O(p))\) and \(W^u(\mathcal O(p))\) denote the stable and unstable manifolds of that orbit. Under the conventional definition, two periodic saddles \(p,q\) are homoclinically related when both
\[
W^u(\mathcal O(p))\pitchfork W^s(\mathcal O(q))\ne\varnothing,
\qquad
W^u(\mathcal O(q))\pitchfork W^s(\mathcal O(p))\ne\varnothing,
\]
where \(\pitchfork\) indicates a transverse intersection. The supplied excerpt does not specify the ambient setting or explicitly confirm this convention.

\textbf{Question 11 (Berger, Bedford, 2023).} For every map \(f_{c,b}\) in the intended family and every infinite sequence \((p_n)_{n\ge1}\) of its periodic saddles such that no two distinct terms are homoclinically related, must
\[
\lim_{n\to\infty}\min\bigl\{|\lambda^-(p_n)|,|\lambda^+(p_n)|\bigr\}=0?
\]
Equivalently, must every \(\varepsilon>0\) admit \(N\ge1\) such that
\[
\min\bigl\{|\lambda^-(p_n)|,|\lambda^+(p_n)|\bigr\}<\varepsilon
\qquad\text{for all }n\ge N?
\]
This is a statement template retaining the source notation. A fully specified mathematical question requires the missing family, domains, parameter restrictions, and definitions of \(\lambda^\pm\); in particular, no choice between eigenvalue multipliers and Lyapunov exponents is made here.
\subsection{Short English statement}
The question asks whether every infinite sequence of pairwise homoclinically unrelated periodic saddles of a fixed map \(f_{c,b}\) must have the smaller absolute value of its two quantities \(\lambda^-\) and \(\lambda^+\) tend to zero. The supplied excerpt leaves the map family and these quantities undefined. Its status is recorded as open in the supplied registry; this is not an independently verified assessment.
\subsection{Sources}
\begin{itemize}
\item Local UnsolvedMath registry, permanent ID \(801059\), exact source problem code \texttt{catalogue entry}, title ``Henon Maps: Open Problems --- Question 11.'' The record attributes the problem to Berger, Bedford, with proposed year 2023, and supplies the status ``open.''
\item Berger, Bedford, \emph{Henon Maps: A List of Open Problems} (2023), Question 11. Supplied source URL: \url{https://arxiv.org/abs/2312.03907}. This reference is supplied by the registry and was not independently consulted for this draft.
\item The supplied provenance records access on 2026-10-16 and extraction by \texttt{arxiv\_latex}. This is supplied dated provenance, not independently verified source access or literature assessment. The record identifies upstream revision \texttt{d371ee917ac9647fbb06cd6fc529deec6f3dda82}.
\end{itemize}
\end{document}
